\documentclass[11pt]{article}

\usepackage[margin=1in]{geometry}
\usepackage{amsmath, amssymb, amsfonts}
\usepackage{booktabs}
\usepackage{enumitem}
\usepackage[colorlinks=true, linkcolor=blue, urlcolor=blue, citecolor=blue]{hyperref}

\title{Notes: Measuring Reinforcement Learning Against the Exact LQR Optimum}
\author{Jon Legasa}
\date{\today}

\begin{document}

\maketitle

\section*{Summary}

LQR has become a standard for measuring RL algorithms, but we found that some algorithms have not
been tested on it, and we are trying to see the impact of LQR across many seeds to see the behavior
of these algorithms. We seek to understand whether there is any way to measure the performance of
these algorithms with respect to each other, and then to measure their performance on RL
environments to see if these differences in performance hold even for many seeds.

%----------------------------------------------------------------------------------------------------
\section{Motivation}

We came together to try to find what makes these algorithms special, following consistent
benchmarks and baselines. We set up the LQR step to motivate the baseline (Section~\ref{sec:riccati}).
Agyei, Sarhadi and Polani (2025) introduced LQI as the baseline for comparing deep RL algorithms on
tracking problems in continuous time, one of them nonlinear (Section~\ref{sec:lqi}). From this
baseline we extend to those setups.

We then introduce the many seeds, the expected value and the average of the learning curves
(Section~\ref{sec:seeds}), because we expect the learning curves to converge to the same value
across seeds, and to be bounded. Think of the set of seeds as points on a line, any of which can be
picked. As we observe all seeds and measure the performance of every algorithm, it could be possible
that, as training time increases, for all seeds on this line the algorithms generate learning curves
that satisfy an ordering inequality. That is the setup of these notes.

The LQR is a standard diagnostic for RL because its optimal controller is known (Recht 2019).
Earlier LQR comparisons covered certainty equivalence, least-squares policy iteration, REINFORCE and
random search (Recht 2019; Mania, Guy and Recht 2018). Agyei et al.\ compared deep RL methods with an
LQI controller using step-response metrics.

%----------------------------------------------------------------------------------------------------
\section{Problem and notation}
\label{sec:problem}

\[
  x_{t+1} = A x_t + B u_t, \qquad c(x, u) = x^\top Q x + u^\top R u, \qquad
  J(\pi) = \sum_{t=0}^{\infty} \gamma^t c(x_t, u_t), \qquad 0 < \gamma \le 1 .
\]

Here $x_t$ is the state at time $t$ and $u_t$ is the control; $A$ describes how the state evolves on
its own and $B$ how the control acts on it. The \textbf{stage cost} is the cost paid at time $t$:
\[
  c(x_t, u_t) = \underbrace{x_t^\top Q\, x_t}_{\text{distance from } 0}
  + \underbrace{u_t^\top R\, u_t}_{\text{control effort}},
  \qquad Q \succeq 0, \quad R \succ 0 .
\]
In the scalar case, $A$, $B$, $Q$ and $R$ are the numbers $a$, $b$, $q$ and $r$.

$J(\pi)$ is the total discounted cost of following the policy $\pi$ from $x_0$: run the system under
$\pi$ and add up the stage costs, with the cost at time $t$ weighted by $\gamma^t$. The discount factor
$\gamma$ is the same one the agents use, so the baseline and the agents optimize the same objective;
$\gamma = 1$ is the undiscounted problem, and every result below reduces to it. The system is
noise-free, with fixed $x_0$ (Section~\ref{sec:limitations} covers noise). Policies are linear,
$u = -Kx$, and $J^* = \min_\pi J(\pi)$ is attained at $K^*$, as derived in Section~\ref{sec:riccati};
both depend on $\gamma$. The agent receives the reward $r_t = -c(x_t, u_t)$ and maximizes the
discounted return $\sum_t \gamma^t r_t = -J(\pi)$. It does not know $A$, $B$, $Q$, $R$, $K^*$ or
$J^*$; these are used only to evaluate it. Numerical examples use $\gamma = 0.99$. Algorithms $A$, $B$ and $C$ (PPO, SAC and TD3) are trained with many
seeds, and $N$ counts training iterations.

%----------------------------------------------------------------------------------------------------
\section{The Riccati solution}
\label{sec:riccati}

The Riccati equation is the Bellman equation of the LQR, reduced to an equation for a single
matrix. The optimal cost-to-go $V(x)$, the total discounted cost of acting optimally from state $x$,
pays the stage cost now and the discounted cost-to-go from the next state:
\[
  V(x) = \min_u \Big[\, x^\top Q x + u^\top R u + \gamma\, V(A x + B u) \,\Big].
\]

Guess that the cost-to-go is quadratic, with a symmetric positive semidefinite matrix $P$; the guess
is confirmed if it leads to a consistent equation for $P$. Substituting it, the bracket becomes
\[
  V(x) = x^\top P x, \qquad x^\top Q x + u^\top R u + \gamma\, (A x + B u)^\top P (A x + B u).
\]
This is a convex quadratic in the control, because $R$ is positive definite and the term with $P$ is
positive semidefinite. Its gradient with respect to the control is
\[
  \nabla_u = 2 R u + 2 \gamma B^\top P (A x + B u)
  = 2\big(R + \gamma B^\top P B\big) u + 2 \gamma B^\top P A\, x ,
\]
and setting it to zero gives the optimal control:
\[
  u^* = -\big(R + \gamma B^\top P B\big)^{-1} \gamma B^\top P A\, x = -K^* x .
\]

Substituting the optimal control back, both sides of the Bellman equation are quadratic forms in
the state:
\[
  x^\top P x = x^\top \Big[\, Q + K^{*\top} R K^* + \gamma (A - B K^*)^\top P (A - B K^*) \,\Big] x .
\]
Expanding the last term and grouping the two terms that contain the gain twice:
\[
  Q + \gamma A^\top P A - \gamma A^\top P B K^* - \gamma K^{*\top} B^\top P A
  + K^{*\top} \big(R + \gamma B^\top P B\big) K^* .
\]
The definition of the optimal gain gives the identity below, which makes the last term equal to the
one before it, so those two cancel:
\[
  \big(R + \gamma B^\top P B\big) K^* = \gamma B^\top P A
  \quad\Longrightarrow\quad
  K^{*\top} \big(R + \gamma B^\top P B\big) K^* = \gamma K^{*\top} B^\top P A .
\]
What remains is the matrix in brackets, and it must equal $P$ because the two quadratic forms agree
for every state and both matrices are symmetric. This is the discounted discrete algebraic Riccati
equation, shown with the optimal gain and the optimal cost:
\begin{gather*}
  P = Q + \gamma A^\top P A - \gamma^2 A^\top P B \big(R + \gamma B^\top P B\big)^{-1} B^\top P A, \\
  K^* = \big(R + \gamma B^\top P B\big)^{-1} \gamma B^\top P A, \qquad J^* = x_0^\top P x_0 .
\end{gather*}

$P$ is the cost-to-go matrix: the optimal cost is the cost-to-go at the initial state, which is why
the optimal cost is the quadratic form of $P$ at the initial state. The equation is algebraic
because, over an infinite horizon, $P$ does not change with time; it is a fixed point of the Bellman
equation. It is quadratic in $P$.

\paragraph{The discount as rescaled dynamics.} Writing $\gamma = \sqrt{\gamma}\cdot\sqrt{\gamma}$ in
every term shows that the discounted equation is the undiscounted one for rescaled matrices:
\[
  A_\gamma = \sqrt{\gamma}\, A, \qquad B_\gamma = \sqrt{\gamma}\, B
  \quad\Longrightarrow\quad
  P = Q + A_\gamma^\top P A_\gamma - A_\gamma^\top P B_\gamma \big(R + B_\gamma^\top P B_\gamma\big)^{-1}
  B_\gamma^\top P A_\gamma .
\]
So a standard Riccati solver applied to $(\sqrt{\gamma} A, \sqrt{\gamma} B, Q, R)$ returns the
discounted $P$, and the same substitution handles the augmented system of Section~\ref{sec:lqi}. The
discounted optimal gain makes $\sqrt{\gamma}(A - BK^*)$ stable, which is weaker than making $A - BK^*$
stable; in the scalar example below it does both.

\paragraph{The scalar case.} With $n = m = 1$:
\[
  A = [a], \quad B = [b], \quad Q = [q], \quad R = [r], \quad P = [p].
\]
Substituting term by term into the Riccati equation, with transposes acting trivially on numbers:
\[
  P = Q + \underbrace{\gamma A^\top P A}_{(1)} - \underbrace{\gamma A^\top P B}_{(2)}\,
  \underbrace{\big(R + \gamma B^\top P B\big)^{-1}}_{(3)}\, \underbrace{\gamma B^\top P A}_{(4)},
\]
\begin{align*}
  (1)\;& \gamma A^\top P A = \gamma\, a\, p\, a = \gamma a^2 p, &
  (2)\;& \gamma A^\top P B = \gamma\, a\, p\, b = \gamma a b\, p, \\
  (3)\;& \big(R + \gamma B^\top P B\big)^{-1} = \frac{1}{r + \gamma b^2 p}, &
  (4)\;& \gamma B^\top P A = \gamma\, b\, p\, a = \gamma a b\, p .
\end{align*}
Multiplying (2), (3) and (4) gives the scalar Riccati equation:
\[
  (2)(3)(4) = \frac{\gamma^2 a^2 b^2 p^2}{r + \gamma b^2 p},
  \qquad
  p = q + \gamma a^2 p - \frac{\gamma^2 a^2 b^2 p^2}{r + \gamma b^2 p} .
\]
With $a = b = q = r = 1$ and $x_0 = 1$, multiply through by $1 + \gamma p$:
\[
  p = 1 + \gamma p - \frac{\gamma^2 p^2}{1 + \gamma p}
  \ \Longrightarrow\ (p - 1 - \gamma p)(1 + \gamma p) = -\gamma^2 p^2
  \ \Longrightarrow\ \gamma p^2 + (1 - 2\gamma)\, p - 1 = 0 ,
\]
and take the root with $p \ge 0$:
\[
  p = \frac{(2\gamma - 1) + \sqrt{4\gamma^2 + 1}}{2\gamma}, \qquad
  K^* = \frac{\gamma a b\, p}{r + \gamma b^2 p} = \frac{\gamma p}{1 + \gamma p}
  = \frac{\sqrt{1 + 4\gamma^2} - 1}{2\gamma}, \qquad
  J^* = p\, x_0^2 = p .
\]
The closed form of the gain follows by substituting $p = K^*/(\gamma(1 - K^*))$, the inverse of the
middle expression, into the quadratic, which leaves
\[
  \gamma (K^*)^2 + K^* - \gamma = 0 .
\]
At $\gamma = 1$ these are the golden-ratio values $p = (1 + \sqrt{5})/2 \approx 1.618$ and
$K^* = (\sqrt{5} - 1)/2 \approx 0.618$. At $\gamma = 0.99$:
\[
  K^* \approx 0.6153, \qquad J^* = p \approx 1.6153 .
\]

So the cost-to-go matrix is the $1 \times 1$ matrix $[p]$, and $V(x) = p x^2$. Two identities are used
below:
\[
  (K^*)^2 = 1 - \frac{K^*}{\gamma}, \qquad 1 + \gamma J^* = 1 + \gamma p = \frac{1}{1 - K^*} .
\]
At $\gamma = 1$ the first reads $(K^*)^2 + K^* - 1 = 0$, and together they give $K^* J^* = 1$. The
gain lies in $(0, 1)$ for every $\gamma \in (0, 1]$, so $1 - K^*$ lies in $(0, 1)$ and the closed loop
$x_{t+1} = (1 - K^*)x_t$ is stable without the discount. In higher dimensions $P$ stays a genuine
matrix, and the Riccati equation is a system of coupled quadratic equations in its entries, usually
solved numerically.

%----------------------------------------------------------------------------------------------------
\section{From LQR to LQI}
\label{sec:lqi}

LQR drives the state to zero. A tracking problem asks an output to follow a reference instead, and
with a constant disturbance LQR leaves a permanent error. LQI is LQR applied to a system enlarged by
one state that accumulates the error.

\paragraph{The offset under a constant disturbance.} Add a constant disturbance to the system and
apply the LQR controller:
\[
  x_{t+1} = A x_t + B u_t + d, \qquad u_t = -K x_t .
\]
At steady state the state stops changing, so it must satisfy
\[
  x_{\mathrm{ss}} = (A - B K)\, x_{\mathrm{ss}} + d
  \quad\Longrightarrow\quad
  x_{\mathrm{ss}} = \big(I - A + B K\big)^{-1} d \neq 0 .
\]
In the scalar example the steady state is a multiple of the disturbance:
\[
  x_{\mathrm{ss}} = (1 - k)\, x_{\mathrm{ss}} + d
  \quad\Longrightarrow\quad k\, x_{\mathrm{ss}} = d
  \quad\Longrightarrow\quad x_{\mathrm{ss}} = \frac{d}{k} \approx 1.625\, d
  \quad (k = K^*,\ \gamma = 0.99).
\]
The controller pushes back only in proportion to the current error, so it needs a nonzero error to
produce the push that cancels the disturbance.

\paragraph{Tracking as shifted regulation.} To make an output follow a constant reference, find the
equilibrium that gives that output, a state and a constant control that hold it:
\[
  y_t = C x_t, \qquad x_{\mathrm{ss}} = A x_{\mathrm{ss}} + B u_{\mathrm{ss}} + d, \qquad
  C x_{\mathrm{ss}} = y_{\mathrm{ref}} .
\]
Subtracting the equilibrium equation from the dynamics removes the disturbance and leaves a pure LQR
problem in shifted coordinates:
\begin{gather*}
  \tilde x_t = x_t - x_{\mathrm{ss}}, \quad \tilde u_t = u_t - u_{\mathrm{ss}}
  \quad\Longrightarrow\quad
  \tilde x_{t+1} = A \tilde x_t + B \tilde u_t, \\
  u_t = u_{\mathrm{ss}} - K\,(x_t - x_{\mathrm{ss}}) .
\end{gather*}
So tracking is regulation around a shifted equilibrium, but computing that equilibrium requires
knowing the disturbance and the model exactly. Without the disturbance the shift is wrong and the
offset returns.

\paragraph{The integral state.} Instead of computing the shift, add a state that accumulates the
tracking error:
\[
  \epsilon_t = y_{\mathrm{ref}} - C x_t, \qquad z_{t+1} = z_t + \epsilon_t
  \quad\Longrightarrow\quad z_t = \sum_{s=0}^{t-1} \epsilon_s .
\]
Here time is discrete, so $z$ is a sum of past errors. In continuous time, where $t$ is a real
number, the same state is an integral, which is where the name comes from:
\[
  \dot z(t) = \epsilon(t) \quad\Longrightarrow\quad z(t) = \int_0^t \epsilon(s)\, ds .
\]

\paragraph{The augmented LQR problem.} Treat $z$ as an extra state component. The enlarged system is
still linear and the discounted cost is still quadratic, with one more weight:
\[
  \begin{pmatrix} x_{t+1} \\ z_{t+1} \end{pmatrix}
  = \begin{pmatrix} A & 0 \\ -C & I \end{pmatrix}
    \begin{pmatrix} x_t \\ z_t \end{pmatrix}
  + \begin{pmatrix} B \\ 0 \end{pmatrix} u_t
  + \begin{pmatrix} d \\ y_{\mathrm{ref}} \end{pmatrix},
  \qquad
  \sum_t \gamma^t \Big( x_t^\top Q x_t + z_t^\top Q_z z_t + u_t^\top R u_t \Big).
\]
The discounted Riccati equation of Section~\ref{sec:riccati}, applied to these larger matrices in the
shifted coordinates (equivalently, the undiscounted one applied to $\sqrt{\gamma}$ times them), gives a
gain with one part for each state, so the LQI controller is proportional plus integral:
\[
  u_t = -K_x x_t - K_z z_t .
\]

\paragraph{Zero steady-state error.} At steady state every state stops changing. For the integral
state this forces the error to zero:
\[
  z_{\mathrm{ss}} = z_{\mathrm{ss}} + \big(y_{\mathrm{ref}} - C x_{\mathrm{ss}}\big)
  \quad\Longrightarrow\quad C x_{\mathrm{ss}} = y_{\mathrm{ref}} .
\]
In the scalar example, with output equal to the state, the two steady-state equations give
\[
  x_{\mathrm{ss}} = y_{\mathrm{ref}}, \qquad
  x_{\mathrm{ss}} = x_{\mathrm{ss}} + u_{\mathrm{ss}} + d
  \ \Longrightarrow\ u_{\mathrm{ss}} = -d = -K_x y_{\mathrm{ref}} - K_z z_{\mathrm{ss}}
  \ \Longrightarrow\ z_{\mathrm{ss}} = \frac{d - K_x y_{\mathrm{ref}}}{K_z} .
\]
The integral state settles at exactly the value that cancels the disturbance, without the
disturbance ever being known. The proportional part no longer needs a nonzero error, because the
accumulated error holds the push.

\paragraph{Continuous time.} The same construction holds with the dynamics written as differential
equations. A continuous discount $e^{-\alpha t}$ is absorbed by the substitution
$\tilde x = e^{-\alpha t/2} x$, $\tilde u = e^{-\alpha t/2} u$, which shifts the dynamics matrix, and
the gain then comes from the continuous algebraic Riccati equation for the shifted matrix:
\begin{gather*}
  \dot x = A x + B u + d, \qquad \dot z = y_{\mathrm{ref}} - C x, \qquad
  \int_0^\infty e^{-\alpha t}\big(x^\top Q x + u^\top R u\big)\, dt, \\
  A_\alpha = A - \tfrac{\alpha}{2} I, \qquad
  A_\alpha^\top P + P A_\alpha - P B R^{-1} B^\top P + Q = 0, \qquad K^* = R^{-1} B^\top P .
\end{gather*}
The agents act at a sampling interval $\Delta t$, so they see a discrete-time system. Holding the
control constant between samples makes the sampled system exact and linear, and the two discounts
match:
\[
  A_d = e^{A \Delta t}, \qquad B_d = \int_0^{\Delta t} e^{A s}\, ds\; B, \qquad
  \gamma = e^{-\alpha \Delta t} .
\]
The agent's reward at each sample defines the discrete stage cost, for example
$\Delta t\,(x_t^\top Q x_t + u_t^\top R u_t)$. The discrete-time baseline of
Section~\ref{sec:riccati}, applied to $(A_d, B_d)$ with that same stage cost and $\gamma$, is the exact
optimum for the problem the agents actually solve, so continuous-time environments need no separate
theory. It differs slightly from the continuous-time optimum above, by an amount that vanishes as the
sampling interval shrinks.

\paragraph{The Agyei et al.\ environments.} Agyei, Sarhadi and Polani (2025) chose LQI because their
tasks are setpoint tracking; in their words, ``the integral action ensures zero steady-state error for
step references, making LQI superior to LQR for tracking applications.'' Their deep RL agents were
trained on the negative of the LQI cost. All four environments are continuous-time tracking
problems, and only one is nonlinear:

\begin{center}
  \begin{tabular}{p{0.3\textwidth} p{0.32\textwidth} p{0.28\textwidth}}
    \toprule
    Environment & Model & Task \\
    \midrule
    Chemical reactor (CSTR) & Linear transfer function, non-minimum phase & Step setpoint tracking \\
    Two-mass spring (ACC 1992 benchmark) & Linear state space & Step setpoint tracking \\
    REMUS underwater vehicle (yaw) & Nonlinear & Step setpoint tracking \\
    Crazyflie quadrotor (altitude) & Reduced-order linear model of a nonlinear system
      & Multi-step altitude tracking, simulation and hardware \\
    \bottomrule
  \end{tabular}
\end{center}

These notes start from the regulation baseline of Sections~\ref{sec:problem} and~\ref{sec:riccati},
where the optimum is exact, and the LQI construction above is how that baseline extends to these
tracking environments.

%----------------------------------------------------------------------------------------------------
\section{The learning curve of one run}
\label{sec:onerun}

\paragraph{Cost of a linear policy.} For the policy with gain $k$, with all parameters equal to one
and initial state one, the trajectory and the stage costs are
\begin{gather*}
  u_t = -k\, x_t, \qquad x_{t+1} = x_t + u_t = (1 - k)\, x_t \ \Longrightarrow\ x_t = (1 - k)^t, \\
  c(x_t, u_t) = x_t^2 + u_t^2 = (1 + k^2)\, x_t^2 = (1 + k^2)(1 - k)^{2t}.
\end{gather*}
The discounted stage costs form a geometric series with ratio $\gamma(1 - k)^2$. It converges exactly
when the ratio is below one:
\begin{gather*}
  \sum_{t=0}^{\infty} \rho^t = \frac{1}{1 - \rho} \ \ (|\rho| < 1), \qquad
  \rho = \gamma (1 - k)^2 < 1 \iff |1 - k| < \frac{1}{\sqrt{\gamma}} \\
  \iff k \in \mathcal{K}_\gamma = \Big(1 - \frac{1}{\sqrt{\gamma}},\ 1 + \frac{1}{\sqrt{\gamma}}\Big).
\end{gather*}
So the cost of a gain is
\[
  J : \mathbb{R} \to [0, \infty], \qquad
  J(k) = \sum_{t=0}^{\infty} \gamma^t (1 + k^2)(1 - k)^{2t}
  = \begin{cases} \dfrac{1 + k^2}{1 - \gamma (1 - k)^2}, & k \in \mathcal{K}_\gamma, \\[1ex]
    +\infty, & \text{otherwise.} \end{cases}
\]
At $\gamma = 1$ the interval is $(0, 2)$ and the denominator is $1 - (1 - k)^2 = k(2 - k)$.

\emph{Finite cost is not stability.} The closed loop is stable when $|1 - k| < 1$, that is
$0 < k < 2$. For $\gamma < 1$ the interval of finite cost is wider: at $\gamma = 0.99$ it is about
$(-0.005, 2.005)$. A gain in the thin bands between the two intervals makes the state grow, yet the
discount shrinks the costs faster than the state grows, so the cost is finite. Such runs are counted
as unstable separately, by the test $|1 - k| \ge 1$, and are not detected by the gap alone.

\paragraph{Functions.} The gap of a gain, and, for one run of an algorithm, its gain, learning curve
and perturbation, for $X \in \{A, B, C\}$, with the convention that a finite number divided by
infinity is zero:
\begin{align*}
  g &: \mathbb{R} \to [0, 1], & g(k) &= \frac{J(k) - J^*}{J(k)} = 1 - \frac{J^*}{J(k)}, \\
  k_X &: \mathbb{N} \to \mathbb{R}, & k_X(N) &= \text{gain after } N \text{ iterations}, \\
  g_X &: \mathbb{N} \to [0, 1], & g_X(N) &= g\big(k_X(N)\big), \\
  e_X &: \mathbb{N} \to \mathbb{R}, & e_X(N) &= k_X(N) - K^* .
\end{align*}
The gap is the excess cost as a fraction of the cost the algorithm actually pays with its gain. Only
the gain changes during training; the optimal cost is fixed.

\emph{Why the gap is bounded.} A gain with finite cost costs at least the optimum, and any other gain
costs infinitely much:
\begin{gather*}
  k \in \mathcal{K}_\gamma: \ J^* \le J(k) < \infty \ \Longrightarrow\ 0 < \frac{J^*}{J(k)} \le 1
  \ \Longrightarrow\ 0 \le g(k) < 1, \\
  k \notin \mathcal{K}_\gamma: \ J(k) = \infty \ \Longrightarrow\ \frac{J^*}{J(k)} = 0
  \ \Longrightarrow\ g(k) = 1 .
\end{gather*}
A run whose gain leaves the finite-cost interval scores one, the worst possible value, instead of
infinity, so the learning curve never blows up.

\emph{An exact formula in this instance.} Putting the excess cost over a common denominator:
Write $h(k) = 1 - \gamma(1 - k)^2$ for the denominator of the cost:
\begin{gather*}
  J(k) - J^* = \frac{1 + k^2}{h(k)} - J^* = \frac{1 + k^2 - J^* + \gamma J^* (1 - k)^2}{h(k)}
  = \frac{(1 + \gamma J^*)\, k^2 - 2\gamma J^* k + (1 - J^* + \gamma J^*)}{h(k)} .
\end{gather*}
The numerator is a quadratic in the gain. A quarter of its discriminant vanishes, because the optimal
cost $p$ satisfies the scalar Riccati equation $\gamma p^2 + (1 - 2\gamma)p - 1 = 0$ of
Section~\ref{sec:riccati}:
\begin{gather*}
  (1 + \gamma J^*)(1 - J^* + \gamma J^*) = 1 + (2\gamma - 1) J^* - \gamma (J^*)^2 + \gamma^2 (J^*)^2, \\
  \gamma^2 (J^*)^2 - (1 + \gamma J^*)(1 - J^* + \gamma J^*)
  = \gamma (J^*)^2 + (1 - 2\gamma) J^* - 1 = 0 .
\end{gather*}
So the numerator is a perfect square, with its double root at the optimal gain:
\[
  k = \frac{2\gamma J^*}{2(1 + \gamma J^*)} = \frac{\gamma p}{1 + \gamma p} = K^*
  \quad\Longrightarrow\quad
  (1 + \gamma J^*)\, k^2 - 2\gamma J^* k + (1 - J^* + \gamma J^*) = (1 + \gamma J^*)(k - K^*)^2 .
\]
Dividing the excess cost by the cost, the denominators cancel:
\begin{gather*}
  g(k) = \frac{(1 + \gamma J^*)(k - K^*)^2}{h(k)} \cdot \frac{h(k)}{1 + k^2}
  = \frac{(1 + \gamma J^*)(k - K^*)^2}{1 + k^2}, \qquad k \in \mathcal{K}_\gamma, \\
  1 + \gamma J^* = \frac{1}{1 - K^*} \approx 2.599 \ \ (\gamma = 0.99), \qquad
  \frac{3 + \sqrt{5}}{2} \approx 2.618 \ \ (\gamma = 1) .
\end{gather*}
This holds only for these parameter values ($a = b = q = r = x_0 = 1$, any $\gamma$); in general only
the definition holds.

\paragraph{Shape of the gap.}
\begin{gather*}
  g(K^*) = 0, \qquad g \ \text{decreasing on } \big(1 - \tfrac{1}{\sqrt{\gamma}}, K^*\big), \quad
  \text{increasing on } \big(K^*, 1 + \tfrac{1}{\sqrt{\gamma}}\big), \\
  g(k) \to 1 \ \text{at both ends of } \mathcal{K}_\gamma, \qquad
  g = 1 \ \text{outside } \mathcal{K}_\gamma .
\end{gather*}
These follow from the derivative of the cost. Differentiating the definition, the constant optimal
cost stays outside:
\[
  g = 1 - J^* J^{-1} \quad\Longrightarrow\quad g' = J^* J^{-2} J' = \frac{J^*}{J^2}\, J' ,
\]
so the gap has the sign of the derivative of the cost. Write the cost as a ratio and differentiate:
\begin{gather*}
  J = \frac{f}{h}, \quad f(k) = 1 + k^2, \quad h(k) = 1 - \gamma (1 - k)^2, \qquad
  J' = \frac{f' h - f h'}{h^2}, \\
  f' h - f h' = 2k \big(1 - \gamma (1 - k)^2\big) - (1 + k^2)\, 2\gamma (1 - k)
  = 2\,(\gamma k^2 + k - \gamma).
\end{gather*}
Since the denominator is positive, the derivative has the sign of the quadratic factor. Its roots
have product $-1$, so exactly one is positive, and it lies in $\mathcal{K}_\gamma$:
\[
  \gamma k^2 + k - \gamma = 0 \ \Longrightarrow\ k = \frac{\sqrt{1 + 4\gamma^2} - 1}{2\gamma} = K^* ,
\]
which recovers the Riccati gain of Section~\ref{sec:riccati}. The other root is below
$-1/\gamma \le 1 - 1/\sqrt{\gamma}$, outside the interval. The derivative is negative to the left of
the optimal gain and positive to the right. At the walls the cost goes to infinity, so the ratio of
the optimal cost to the cost goes to zero and the gap goes to one. Together with the value one outside
the finite-cost interval, this makes the gap continuous on the whole real line.

Geometrically, the gap is a valley with its bottom at the optimal gain, rising to height one at the
walls $1 \pm 1/\sqrt{\gamma}$ and staying there beyond them. The perturbation is the horizontal
distance from the bottom, and the gap is the height. The left wall is closer to the bottom than the
right wall (about $0.62$ against $1.39$ at $\gamma = 0.99$), so the valley is steeper on the left.

\paragraph{Gap and perturbation.} Two facts connect the size of the gap to the size of the
perturbation.

\emph{A small perturbation gives a small gap.} The gap is continuous, as shown above. Continuity at
the optimal gain, where the gap is zero, says exactly:
\[
  \forall \varepsilon > 0 \ \exists \delta > 0 : \quad
  |k - K^*| < \delta \ \Longrightarrow\ g(k) < \varepsilon .
\]

\emph{A small gap gives a small perturbation.} This holds because the optimal gain is the unique
minimizer:
\[
  \forall \eta > 0 \ \exists \varepsilon > 0 : \quad
  g(k) < \varepsilon \ \Longrightarrow\ |k - K^*| < \eta .
\]
To find such an $\varepsilon$, take the smaller of the two gaps at distance $\eta$ from the optimal
gain, for $\eta$ smaller than the distance $K^* - (1 - 1/\sqrt{\gamma})$ to the left wall:
\[
  \varepsilon = \min\big\{\, g(K^* - \eta),\ g(K^* + \eta) \,\big\} \in (0, 1) .
\]
The gap decreases toward the optimal gain from the left and increases away from it on the right, and
every gain outside the finite-cost interval has gap one, which is larger than $\varepsilon$. So every gain
at distance $\eta$ or more has a gap of at least $\varepsilon$, and a gap below $\varepsilon$ forces a
distance below $\eta$.

Applying the two facts along the sequence of gains of a run: once the learning curve stays below
$\varepsilon$, the perturbation stays below $\eta$, and once the perturbation stays below $\delta$,
the learning curve stays below $\varepsilon$. This gives
\[
  g_X(N) \to 0 \iff k_X(N) \to K^* \iff e_X(N) \to 0 .
\]
Both facts hold for any continuous gap with a unique minimum, not only for this example.

\paragraph{Behaviour near the optimal gain.} Dividing the exact formula by the square of the
perturbation:
\[
  \frac{g(K^* + e)}{e^2} = \frac{1 + \gamma J^*}{1 + (K^* + e)^2}
  \ \xrightarrow{\ e \to 0\ }\ c_\gamma = \frac{1 + \gamma J^*}{1 + (K^*)^2} .
\]
Using the two identities of Section~\ref{sec:riccati} to simplify the limit:
\[
  1 + \gamma J^* = \frac{1}{1 - K^*}, \qquad 1 + (K^*)^2 = 2 - \frac{K^*}{\gamma}
  \quad\Longrightarrow\quad
  c_\gamma = \frac{1}{(1 - K^*)\big(1 + (K^*)^2\big)} .
\]
At $\gamma = 1$ this is $(3 + \sqrt{5})/(5 - \sqrt{5}) = 1 + 2/\sqrt{5} \approx 1.894$; at
$\gamma = 0.99$ it is $c_{0.99} \approx 1.885$. So near the optimum the gap is approximately $c_\gamma$
times the squared perturbation. For a perturbation of at most 0.1 the gain is positive, so the
denominator increases with the perturbation and the ratio decreases; its extremes are at the two ends,
which gives exact bounds, shown at $\gamma = 0.99$ (rounded outward):
\begin{gather*}
  |e| \le 0.1: \quad \frac{1 + \gamma J^*}{1 + (K^* + 0.1)^2} \le \frac{g(K^* + e)}{e^2}
  \le \frac{1 + \gamma J^*}{1 + (K^* - 0.1)^2} \\
  \Longrightarrow\quad 1.719\, e^2 \le g(K^* + e) \le 2.054\, e^2 .
\end{gather*}
The bounds give explicit choices in the two facts. For the first, the upper bound controls the gap
by the perturbation; for the second, the lower bound and the monotonicity above control the
perturbation by the gap:
\[
  |e| < \delta = \sqrt{\varepsilon / 2.054} \le 0.1
  \ \Longrightarrow\ g(K^* + e) \le 2.054\, e^2 < 2.054\, \delta^2 = \varepsilon ,
\]
\[
  \eta \le 0.1, \ \varepsilon = 1.719\, \eta^2: \quad
  |e| \ge \eta \ \Longrightarrow\ g(K^* + e) \ge \min\big\{ g(K^* - \eta), g(K^* + \eta) \big\}
  \ge 1.719\, \eta^2 = \varepsilon .
\]
The two observations below come from the exact formula at $\gamma = 0.99$. The two gains are the same
distance from the optimum, so the numerators agree and only the denominators differ:
\begin{gather*}
  K^* - 0.118 \approx 0.497, \qquad K^* + 0.118 \approx 0.733, \\
  g(0.497) = \frac{2.599 \cdot 0.118^2}{1 + 0.497^2} \approx 0.029, \qquad
  g(0.733) = \frac{2.599 \cdot 0.118^2}{1 + 0.733^2} \approx 0.024 .
\end{gather*}
\begin{itemize}[leftmargin=*]
  \item \textbf{Observation 1.} A gain about 19\% below the optimum has a gap of only about 3\%: the
    gap is quadratic in the perturbation.
  \item \textbf{Observation 2.} Undershooting and overshooting the optimal gain by the same amount,
    0.118, give gaps of 0.029 and 0.024: undershooting costs more, because the smaller gain has the
    smaller denominator.
\end{itemize}
Both observations hold at every $\gamma$ near one; the numbers change only in the third decimal.

\paragraph{Limits.} The gap is continuous on the whole real line, so continuity carries any limit of
the gains through the gap, and the floor of the learning curve is the gap at the limit gain:
\[
  k_X(N) \to k_{X,\infty} \in \mathbb{R} \ \Longrightarrow\ g_X(N) \to g\big(k_{X,\infty}\big) \in [0, 1] .
\]
A run that settles at a gain outside the finite-cost interval has floor one; one that settles in the
thin unstable bands inside it has a floor below one, and is flagged by the stability test.

\paragraph{Ordering.} The learning curves of the three algorithms may be ordered at every iteration,
or only after some iteration:
\[
  g_A(N) \le g_B(N) \le g_C(N) \quad \forall N, \qquad
  \exists\, N_0 : \ g_A(N) \le g_B(N) \le g_C(N) \quad \forall N \ge N_0 .
\]
The smallest such iteration is the crossing time. Near the optimal gain the gap is, to leading
order, an increasing function of the size of the perturbation, so the ordering of the curves follows
the ordering of the perturbations. If the perturbations decay like power laws, the ratio of two of
them shows which one wins eventually and when they cross:
\[
  |e_X(N)| \approx \kappa_X N^{-\beta_X}
  \ \Longrightarrow\
  \frac{|e_X(N)|}{|e_Y(N)|} \approx \frac{\kappa_X}{\kappa_Y}\, N^{-(\beta_X - \beta_Y)} \to 0
  \quad \text{if } \beta_X > \beta_Y ,
\]
\[
  \frac{\kappa_X}{\kappa_Y}\, N^{-(\beta_X - \beta_Y)} = 1
  \iff N = N_{XY} = \left(\frac{\kappa_X}{\kappa_Y}\right)^{1/(\beta_X - \beta_Y)} .
\]
So the faster-decaying algorithm is eventually below; if it starts above, it crosses below at the
iteration shown.

%----------------------------------------------------------------------------------------------------
\section{Many seeds}
\label{sec:seeds}

A seed is a random number. All the randomness in one training run, the initial parameters, the
exploration noise and any noise in the environment, is generated from it, so once the seed is picked
the run is deterministic. The set of seeds $\Omega$ is a set of random numbers; in code it is a
finite set of integers, which can be pictured as points on a line:
\[
  \Omega = \{\, 0, 1, 2, \dots, 2^{32} - 1 \,\} \subset \mathbb{R}, \qquad \omega \in \Omega .
\]
Picking a seed is drawing a point from this line, and each point gives one learning curve. Over all
of $\Omega$, an algorithm therefore produces a collection of learning curves, one per seed, and the
analysis of the algorithm is the analysis of this collection: its average, its spread, its bounds and
where it settles. The question of the notes is whether, as training time increases, the ordering of
Section~\ref{sec:onerun} holds for every seed on this line.

The curves are bounded below, because no gain does better than the optimum:
\[
  J(k) \ge J^* \quad\Longrightarrow\quad g_X(\omega, N) \ge 0 \quad \text{for every seed and iteration} .
\]
They are also bounded above by one, by the definition of the gap in Section~\ref{sec:onerun}. So
every expected curve below is finite at every iteration, and runs that leave the finite-cost
interval stay in the averages with the value one:
\[
  0 \le g_X(\omega, N) \le 1 \quad \text{for every seed and iteration} .
\]

\paragraph{Functions.} Seeds are drawn uniformly from the set of integer seeds, and the gain and
learning curve become functions of the seed as well as the iteration:
\[
  \omega \in \Omega, \qquad P(\omega) = 2^{-32},
\]
\begin{align*}
  k_X &: \Omega \times \mathbb{N} \to \mathbb{R}, &
    k_X(\omega, N) &= \text{gain after } N \text{ iterations with seed } \omega, \\
  g_X &: \Omega \times \mathbb{N} \to [0, 1], &
    g_X(\omega, N) &= g\big(k_X(\omega, N)\big).
\end{align*}
Fixing $\omega$ recovers Section~\ref{sec:onerun}.

\paragraph{Expected and estimated curves.} Over all seeds, and over the $M$ seeds run:
\[
  \bar g_X(N) = \sum_{\omega \in \Omega} P(\omega)\, g\big(k_X(\omega, N)\big), \qquad
  \hat g_X(N) = \frac{1}{M} \sum_{j=1}^{M} g\big(k_X(\omega_j, N)\big).
\]
The seeds are independent, so the variance of an average of $M$ curves is the variance of one curve
divided by $M$:
\begin{gather*}
  \operatorname{Var}\big[\hat g_X(N)\big] = \frac{1}{M^2} \sum_{j=1}^{M}
    \operatorname{Var}\big[g_X(\omega_j, N)\big] = \frac{\sigma_X(N)^2}{M}, \\
  \operatorname{SE}\big[\hat g_X(N)\big] = \frac{\sigma_X(N)}{\sqrt{M}}, \qquad
  \sigma_X(N)^2 = \operatorname{Var}\big[g_X(\cdot, N)\big].
\end{gather*}
Runs whose gain has infinite discounted cost stay in the averages with gap one; the fraction of
unstable runs, $|1 - k_X(\omega, N)| \ge 1$, is reported alongside the curve.

\paragraph{Convergence of the estimated curve.} Fix an iteration. Each seed run gives one number
between zero and one, the seeds are drawn independently, and the mean of each number is the expected
curve:
\[
  Y_j = g_X(\omega_j, N) \in [0, 1], \qquad
  \mathbb{E}[Y_j] = \sum_{\omega \in \Omega} P(\omega)\, g_X(\omega, N) = \bar g_X(N), \qquad
  \hat g_X(N) = \frac{1}{M} \sum_{j=1}^{M} Y_j .
\]
Linearity of expectation gives the mean of the estimated curve, and its variance was computed above:
\[
  \mathbb{E}\big[\hat g_X(N)\big] = \frac{1}{M} \sum_{j=1}^{M} \mathbb{E}[Y_j] = \bar g_X(N), \qquad
  \operatorname{Var}\big[\hat g_X(N)\big] = \frac{\sigma_X(N)^2}{M} .
\]
Because the gap lies between zero and one, the variance of one run is at most one quarter, whatever
the algorithm does. A number between zero and one is at least its square, and a downward parabola
with roots at zero and one has its maximum, one quarter, at one half:
\[
  Y_j^2 \le Y_j \ \Longrightarrow\ \sigma_X(N)^2 = \mathbb{E}[Y_j^2] - \bar g_X(N)^2
  \le \bar g_X(N) - \bar g_X(N)^2 = \bar g_X(N)\big(1 - \bar g_X(N)\big) \le \frac{1}{4} .
\]
Applying the inequalities of Section~\ref{sec:ineq} to the estimated curve gives the weak law of large
numbers in two forms:
\begin{align*}
  \text{(Chebyshev)} \quad & P\Big(\big|\hat g_X(N) - \bar g_X(N)\big| \ge \varepsilon\Big)
    \le \frac{\sigma_X(N)^2}{M \varepsilon^2} \le \frac{1}{4 M \varepsilon^2}, \\
  \text{(Hoeffding)} \quad & P\Big(\big|\hat g_X(N) - \bar g_X(N)\big| \ge \varepsilon\Big)
    \le 2 e^{-2 M \varepsilon^2} .
\end{align*}
Both bounds go to zero as the number of seeds grows. Setting the Hoeffding bound below a chosen risk
and solving for the number of seeds:
\[
  2 e^{-2 M \varepsilon^2} \le \delta \iff -2 M \varepsilon^2 \le \ln \frac{\delta}{2}
  \iff M \ge \frac{\ln(2/\delta)}{2 \varepsilon^2} .
\]
In the form of Section~\ref{sec:onerun}, for every tolerance and every risk there is a number of
seeds that achieves both:
\[
  \forall \varepsilon > 0 \ \forall \delta > 0 : \quad
  M \ge \frac{\ln(2/\delta)}{2 \varepsilon^2}
  \ \Longrightarrow\ P\Big(\big|\hat g_X(N) - \bar g_X(N)\big| \ge \varepsilon\Big) \le \delta .
\]
For example, to be within 0.05 of the expected curve with probability 0.95, the two inequalities ask
for
\[
  \varepsilon = \delta = 0.05: \qquad
  \frac{1}{4 \varepsilon^2 \delta} = \frac{1}{4 \cdot 0.0025 \cdot 0.05} = 2000 \ \text{seeds}, \qquad
  \frac{\ln(2/\delta)}{2 \varepsilon^2} = \frac{\ln 40}{0.005} \approx 738 \ \text{seeds} .
\]
The strong law of large numbers, which is cited here and not derived, gives the stronger statement
that the estimated curve converges to the expected curve with probability one:
\[
  P\Big(\lim_{M \to \infty} \hat g_X(N) = \bar g_X(N)\Big) = 1 .
\]
The expected curve is never computed exactly. The statements above say how close the estimate is and
with what probability, which is what a confidence band reports. Each statement holds for one
iteration at a time; for several checkpoints together the risks add, because the probability of a
union is at most the sum of the probabilities:
\[
  P\Big(\exists\, i \le n : \big|\hat g_X(N_i) - \bar g_X(N_i)\big| \ge \varepsilon\Big)
  \le \sum_{i=1}^{n} P\Big(\big|\hat g_X(N_i) - \bar g_X(N_i)\big| \ge \varepsilon\Big)
  \le 2 n\, e^{-2 M \varepsilon^2} .
\]
The expected curve belongs to one algorithm with fixed hyperparameters; changing, for example, the
learning rate changes the expected curve.

\paragraph{Bias and variance.} Averaging the near-optimum approximation of Section~\ref{sec:onerun} over seeds gives
\begin{gather*}
  \bar g_X(N) \approx c_\gamma\, \mathbb{E}\big[e_X(\cdot, N)^2\big], \\
  \mathbb{E}[e^2] = \big(\mathbb{E}[e]\big)^2 + \operatorname{Var}[e], \qquad
  \mathbb{E}[e] = \mathbb{E}[k] - K^*, \qquad \operatorname{Var}[e] = \operatorname{Var}[k],
\end{gather*}
where the middle identity is the definition of variance rearranged, and the last two hold because
the perturbation is the gain shifted by a constant. Together:
\[
  \bar g_X(N) \approx c_\gamma \Big[\big(\mathbb{E}[k_X(\cdot, N)] - K^*\big)^2
  + \operatorname{Var}\big[k_X(\cdot, N)\big]\Big].
\]
\begin{itemize}[leftmargin=*]
  \item \textbf{Observation 3.} The expected gap splits into squared bias (settling at the wrong gain
    on average) and variance (inconsistency across seeds).
\end{itemize}

\paragraph{Ordering across seeds.} From strongest to weakest:
\[
  \max_{\omega} g_A(\omega, N) \le \min_{\omega'} g_B(\omega', N), \qquad
  \bar g_A(N) \le \bar g_B(N), \qquad
  \lim_{N \to \infty} \bar g_A(N) \le \lim_{N \to \infty} \bar g_B(N),
\]
that is, band separation, expected-curve ordering, and floor ordering. Each implies the next. Band
separation places every run of $A$ below every run of $B$, so averaging over both sets of seeds keeps
the inequality; taking limits keeps it again:
\begin{gather*}
  g_A(\omega, N) \le \max_{\omega} g_A(\omega, N) \le \min_{\omega'} g_B(\omega', N)
    \le g_B(\omega', N) \\
  \Longrightarrow\ \bar g_A(N) \le \bar g_B(N)
  \ \Longrightarrow\ \lim_{N \to \infty} \bar g_A(N) \le \lim_{N \to \infty} \bar g_B(N).
\end{gather*}

The converses, the same implications read backwards, fail. Each reversed arrow fails for its own
reason, shown by one counterexample.

\emph{Floor ordering does not imply expected-curve ordering.} A floor depends only on the tail of a
curve, so two curves can end in one order after crossing early on:
\[
  \bar g_A(N) = 0.1 + \frac{0.8}{N}, \qquad \bar g_B(N) = 0.2 + \frac{0.1}{N}, \qquad N \ge 1: \qquad
  \lim_{N \to \infty} \bar g_A = 0.1 < 0.2 = \lim_{N \to \infty} \bar g_B ,
\]
\[
  \bar g_A(N) > \bar g_B(N) \iff 0.1 + \frac{0.8}{N} > 0.2 + \frac{0.1}{N}
  \iff \frac{0.7}{N} > 0.1 \iff N < 7 .
\]
Both curves stay between zero and one, as gaps must. Algorithm $A$ starts much worse and overtakes
$B$ only after 7 iterations; the floors alone hide this.

\emph{Expected-curve ordering does not imply band separation.} An average can be ordered while
individual runs are not. Suppose two seeds of each algorithm give these gaps at some iteration:
\begin{center}
  \begin{tabular}{lccc}
    \toprule
      & Seed 1 & Seed 2 & Average \\
    \midrule
    $A$ & 0.01 & 0.05 & 0.03 \\
    $B$ & 0.02 & 0.06 & 0.04 \\
    \bottomrule
  \end{tabular}
\end{center}
The averages are ordered, 0.03 below 0.04, but the worst run of $A$, 0.05, is above the best run of
$B$, 0.02. The two sets of runs overlap, so not every run of $A$ beats every run of $B$.

For an experiment, this means each claim must be checked on its own. Ordered floors say nothing about
early training, and ordered averages say nothing about individual runs; only band separation supports
a statement about every run at every iteration.

%----------------------------------------------------------------------------------------------------
\section{The inequalities of Chebyshev and Hoeffding}
\label{sec:ineq}

Both inequalities bound the probability that a random quantity lands far from its mean, and both are
derived from Markov's inequality.

\paragraph{Markov's inequality.} For a nonnegative random variable and a positive level, the
indicator of the event is at most the variable divided by the level, and taking expectations gives
the inequality:
\[
  Z \ge 0, \ a > 0: \qquad \mathbf{1}\{Z \ge a\} \le \frac{Z}{a}
  \ \Longrightarrow\ P(Z \ge a) = \mathbb{E}\big[\mathbf{1}\{Z \ge a\}\big] \le \frac{\mathbb{E}[Z]}{a} .
\]
The pointwise inequality holds because, when the event happens, the left side is one and the right
side is at least one, and otherwise the left side is zero and the right side is nonnegative.

\paragraph{Chebyshev's inequality.} Apply Markov's inequality to the squared distance from the mean,
with the squared tolerance as the level:
\[
  P\big(|Y - \mathbb{E}Y| \ge \varepsilon\big) = P\big((Y - \mathbb{E}Y)^2 \ge \varepsilon^2\big)
  \le \frac{\mathbb{E}\big[(Y - \mathbb{E}Y)^2\big]}{\varepsilon^2} = \frac{\operatorname{Var} Y}{\varepsilon^2} .
\]
It needs only a finite variance. Applied to the estimated curve, whose variance is the variance of one
run divided by the number of seeds, it gives the first bound of Section~\ref{sec:seeds}, which decays
only like one over the number of seeds.

\paragraph{Hoeffding's inequality.} For independent draws between zero and one with a common mean:
\[
  Y_1, \dots, Y_M \in [0, 1] \ \text{independent}, \quad \mathbb{E}[Y_j] = m: \qquad
  P\Big(\Big|\frac{1}{M} \sum_{j=1}^{M} Y_j - m\Big| \ge \varepsilon\Big) \le 2 e^{-2 M \varepsilon^2} .
\]
The derivation has three steps.

\emph{Step 1: Markov's inequality on an exponential.} For a positive parameter the exponential is
increasing, so the event can be rewritten and Markov's inequality applied; independence turns the
expectation of the product into the product of the expectations:
\begin{align*}
  s > 0: \quad P\Big(\sum_{j=1}^{M} (Y_j - m) \ge M \varepsilon\Big)
  &= P\Big(e^{s \sum_j (Y_j - m)} \ge e^{s M \varepsilon}\Big)
  \le e^{-s M \varepsilon}\, \mathbb{E}\Big[e^{s \sum_j (Y_j - m)}\Big] \\
  &= e^{-s M \varepsilon} \prod_{j=1}^{M} \mathbb{E}\big[e^{s (Y_j - m)}\big] .
\end{align*}

\emph{Step 2: Hoeffding's lemma.} For one draw between zero and one with mean $m$:
\[
  \mathbb{E}\big[e^{s (Y - m)}\big] \le e^{s^2 / 8} .
\]
The exponential is convex, so on the interval from zero to one it lies below its chord. Taking
expectations bounds the exponential moment by the exponential of a function $L$:
\begin{gather*}
  e^{s y} \le (1 - y) + y\, e^{s} \ \ (0 \le y \le 1)
  \ \Longrightarrow\ \mathbb{E}\big[e^{s Y}\big] \le 1 - m + m e^{s}
  \ \Longrightarrow\ \mathbb{E}\big[e^{s (Y - m)}\big] \le e^{L(s)}, \\
  L(s) = -s m + \ln\big(1 - m + m e^{s}\big) .
\end{gather*}
Differentiating $L$, with $q$ the fraction shown, which lies between zero and one:
\begin{gather*}
  q(s) = \frac{m e^{s}}{1 - m + m e^{s}}, \qquad L'(s) = -m + q(s), \qquad
  q'(s) = q(s)\big(1 - q(s)\big), \\
  L(0) = 0, \qquad L'(0) = 0, \qquad L''(s) = q(s)\big(1 - q(s)\big) \le \frac{1}{4} ,
\end{gather*}
where the bound on the second derivative is the same parabola argument as in
Section~\ref{sec:seeds}. Taylor's theorem with the remainder term then gives the lemma:
\[
  L(s) = L(0) + L'(0)\, s + \frac{L''(\xi)}{2}\, s^2 \le \frac{s^2}{8}
  \quad \text{for some } \xi \text{ between } 0 \text{ and } s .
\]

\emph{Step 3: choose the parameter.} Combining the first two steps and minimizing the exponent over
the parameter:
\begin{gather*}
  P\Big(\frac{1}{M} \sum_{j=1}^{M} Y_j - m \ge \varepsilon\Big) \le e^{-s M \varepsilon + M s^2 / 8}, \\
  \frac{d}{ds}\Big(-s M \varepsilon + \frac{M s^2}{8}\Big) = 0 \iff s = 4 \varepsilon
  \ \Longrightarrow\ e^{-4 M \varepsilon^2 + 2 M \varepsilon^2} = e^{-2 M \varepsilon^2} .
\end{gather*}
The same argument applied to one minus each draw bounds the lower tail by the same amount, and adding
the two tails gives the factor two.

\paragraph{Comparison.} The number of seeds each inequality asks for, with the variance bounded by one
quarter in Chebyshev's case:
\[
  \text{Chebyshev: } \ M \ge \frac{1}{4 \varepsilon^2 \delta}, \qquad
  \text{Hoeffding: } \ M \ge \frac{\ln(2/\delta)}{2 \varepsilon^2} .
\]
\begin{itemize}[leftmargin=*]
  \item Chebyshev needs only a finite variance; Hoeffding needs the draws to be bounded, which the gap
    provides.
  \item Halving the risk doubles Chebyshev's count but adds only a constant to Hoeffding's, because
    the risk enters through a logarithm.
  \item Hoeffding uses only the range of the gap, not its actual spread. When the runs of an algorithm
    are consistent, so the variance is far below one quarter, Chebyshev with an estimated variance, or
    a variance-aware inequality such as Bernstein's, can ask for fewer seeds.
\end{itemize}

%----------------------------------------------------------------------------------------------------
\section{A nonlinear benchmark with an exact solution}
\label{sec:nlqr}

The LQR baseline is exact because its Bellman equation reduces to the Riccati equation. For a
nonlinear system the Bellman equation has no closed-form solution in general, so there is no $J^*$
and no gap. This section builds the simplest nonlinear problem we found in which the solution stays
exact: an LQR problem is bent so that the Riccati matrix of the linear problem is still the
cost-to-go of the nonlinear one. Everything stays in discrete time with the same discount, so $J^*$
is the true optimum of the very problem the agents are run on.

\paragraph{The linear starting point.} Take an LQR problem of Section~\ref{sec:problem} with $n$
states and one control, $R = r$, with two properties. The control enters the last state only, and
the Riccati matrix does not couple the last state to the others:
\[
  B = b\, e_n, \qquad P e_n = p\, e_n ,
\]
where $e_n$ is the last unit vector. A diagonal $P$ is the simplest case, and every scalar problem
qualifies. Write $y = A x$ for the next state of the uncontrolled linear system and $y_n$ for its
last coordinate.

\paragraph{The nonlinear system.} Replace the constant input gain $b$ by a gain that depends on the
first state, and rescale the last row of the drift:
\begin{gather*}
  x^+_i = y_i \quad (i < n), \qquad x^+_n = s(x)\, y_n + \beta(x)\, u, \\
  \beta(x) = b\, \frac{2 + \cos(k x_1)}{3}, \qquad
  s(x) = \sqrt{\frac{r + \gamma p\, \beta(x)^2}{r + \gamma p\, b^2}} .
\end{gather*}
The stage cost $x^\top Q x + r u^2$ and the discount are unchanged. The gain $\beta$ moves between
$b/3$ and $b$ and never vanishes, so the control always acts. At $x_1 = 0$ we have $\beta = b$ and
$s = 1$: the linearization at the origin is the linear problem. The number $k \ge 0$ sets how fast
the gain varies; at $k = 0$ the system is the linear problem exactly.

\paragraph{The cost-to-go is still $x^\top P x$.} Substitute $V(x) = x^\top P x$ into the Bellman
equation of the nonlinear system. Because $P$ does not couple the last state to the others, the
next-state term splits into a part without the control and a square:
\[
  x^{+\top} P x^+ = \sum_{i, j < n} P_{ij}\, y_i y_j + p\, \big(s\, y_n + \beta u\big)^2 .
\]
This is where the structure of $P$ matters: with a cross term between $x_n$ and another state, the
control would also multiply the other coordinates of $y$ and the steps below would not close. The
control appears only in $r u^2 + \gamma p (s y_n + \beta u)^2$, a convex quadratic, minimized at
\[
  u^* = -\frac{\gamma p\, \beta\, s\, y_n}{r + \gamma p\, \beta^2},
  \qquad
  \min_u \Big[ r u^2 + \gamma p\, (s y_n + \beta u)^2 \Big]
  = \frac{\gamma p\, r\, s^2}{r + \gamma p\, \beta^2}\; y_n^2
  = \frac{\gamma p\, r}{r + \gamma p\, b^2}\; y_n^2 .
\]
The last equality is the definition of $s$: the scale is chosen so that the minimized term does not
see the nonlinearity. The minimized bracket of the Bellman equation is therefore
\[
  x^\top Q x + \gamma \sum_{i, j < n} P_{ij}\, y_i y_j + \frac{\gamma p\, r}{r + \gamma p\, b^2}\; y_n^2 ,
\]
which is the minimized bracket of the linear problem ($\beta = b$, $s = 1$), and that equals
$x^\top P x$ by the Riccati equation of Section~\ref{sec:riccati}. So $V(x) = x^\top P x$ satisfies
the Bellman equation of the nonlinear problem, for every $k$.

\paragraph{It is the optimum.} Assume $Q \succ 0$. The Bellman equation says that for every state
and control $V(x) \le c(x, u) + \gamma V(x^+)$, with equality at $u^*$. Applying it along the
trajectory of any policy $\pi$ for $T$ steps,
\[
  V(x_0) \le \sum_{t=0}^{T-1} \gamma^t c(x_t, u_t) + \gamma^T V(x_T) .
\]
If $J(\pi)$ is finite, then $\gamma^t x_t^\top Q x_t \to 0$, hence $\gamma^T V(x_T) \to 0$, and the
limit gives $V(x_0) \le J(\pi)$. Under $u^*$ the inequality is an equality and $V \ge 0$, so the
partial sums of the cost never exceed $V(x_0)$, and $J(u^*) \le V(x_0)$. Together:
\[
  J^* = x_0^\top P x_0, \qquad
  u^*(x) = -\frac{\gamma p\, \beta(x)\, s(x)}{r + \gamma p\, \beta(x)^2}\; (A x)_n .
\]
The optimal cost is the same number as for the linear problem. The optimal policy is not: it agrees
with $u = -K^* x$ where $x_1 = 0$ and differs elsewhere, so no linear gain is optimal when $k > 0$.

\paragraph{The instance.} Two states, $Q = I$, $R = 1$, $b = 1$, and a diagonal Riccati matrix chosen
in advance, $P = \mathrm{diag}(p_1, p_2)$. With $a_1^\top$, $a_2^\top$ the rows of $A$, the Riccati
equation for a diagonal $P$ and $B = b\, e_2$ reads
\[
  \mathrm{diag}(p_1 - 1,\ p_2 - 1) = \gamma p_1\, a_1 a_1^\top
  + \frac{\gamma p_2}{1 + \gamma p_2 b^2}\; a_2 a_2^\top ,
\]
so $A$ is found by choosing $a_{11}$, taking $a_{12} > 0$ from
$a_{11}^2/(p_1 - 1) + a_{12}^2/(p_2 - 1) = 1/(\gamma p_1)$, which makes the left side minus the first
term a matrix of rank one, and reading $a_2$ off it. The matrix $A$ depends on $\gamma$. With
$p_1 = 2$, $p_2 = 3/2$, $a_{11} = 1/2$ and $\gamma = 0.99$:
\begin{gather*}
  A = \begin{pmatrix} 0.5 & 0.3571 \\ 0.9193 & -0.6436 \end{pmatrix}, \qquad
  B = \begin{pmatrix} 0 \\ 1 \end{pmatrix}, \qquad
  P = \begin{pmatrix} 2 & 0 \\ 0 & 1.5 \end{pmatrix}, \qquad
  K^* = \begin{pmatrix} 0.5493 & -0.3846 \end{pmatrix}, \\
  x_0 = (1, -1), \qquad J^* = 2 x_{0,1}^2 + 1.5\, x_{0,2}^2 = 3.5 .
\end{gather*}
Written out, the system is
\[
  x_1^+ = 0.5\, x_1 + 0.3571\, x_2, \qquad
  x_2^+ = s(x_1)\, \big(0.9193\, x_1 - 0.6436\, x_2\big) + \beta(x_1)\, u .
\]
The spectral radius of $A$ is $0.88$, so the uncontrolled system is stable. This is deliberate: the
environment never clips the state or ends an episode early, and an untrained agent does not
stabilize the system, so an open-loop unstable instance would diverge before learning starts.

The scalar running example qualifies as well ($n = 1$, $a = b = q = r = 1$): it becomes
$x^+ = s(x)\, x + \beta(x)\, u$ with the same $J^* = p$ as in Section~\ref{sec:riccati}.

\paragraph{The gap.} For any deterministic policy $\pi$ the gap is defined as in
Section~\ref{sec:onerun},
\[
  g(\pi) = 1 - \frac{J^*}{J(\pi)} \in [0, 1],
\]
with the same convention that a policy of infinite cost has gap one. It lies in $[0, 1]$ because
$J^*$ is the true optimum. Four things differ from the linear case.
\begin{itemize}[leftmargin=*]
  \item \textbf{The cost of a policy is measured, not computed.} There is no closed form for the
    cost of a gain and no finite-cost interval $\mathcal{K}_\gamma$; $J(\pi)$ is the discounted cost
    of a long deterministic rollout from $x_0$, and stability is judged from the same rollout.
  \item \textbf{Linear policies have a floor.} The best gain $u = -Kx$ has a gap above zero. An
    algorithm restricted to linear policies converges to this floor at best; a policy class that
    contains $u^*$ (a table, a network) can reach zero. On this problem the floor of a learning
    curve separates what the policy class can represent from what the algorithm manages to learn.
  \item \textbf{The LQR design is suboptimal.} The gain $K^*$ of the linearization, which is what
    LQR applied to the linearized system returns, has a positive gap.
  \item \textbf{No expansion around an optimal gain.} The bias and variance split of
    Section~\ref{sec:seeds} expands the gap around $K^*$ and needs the optimum to be a gain; it does
    not apply here.
\end{itemize}
The table gives these gaps for the instance above; the best linear gain is found by a numerical
search over $K$.
\begin{center}
  \begin{tabular}{@{}cccc@{}}
    \toprule
    $k$ & $g(u = 0)$ & $g(u = -K^* x)$ & $\min_K g(u = -Kx)$ \\
    \midrule
    0 & 0.654 & 0      & 0 \\
    1 & 0.583 & 0.0026 & 0.0002 \\
    2 & 0.416 & 0.0308 & 0.0044 \\
    3 & 0.299 & 0.0703 & 0.0116 \\
    4 & 0.297 & 0.0449 & 0.0037 \\
    \bottomrule
  \end{tabular}
\end{center}
The effect is not monotone in $k$: both gaps are smaller at $k = 4$ than at $k = 3$. We have not
analysed why. The floors are small, about one percent at most, so they are visible only on a
logarithmic gap axis.

\emph{$k$ and the size of $x_0$ are the same dial.} The nonlinearity enters only through $k x_1$,
and the system is otherwise homogeneous of degree one. Replacing $(x, u, k)$ by
$(\lambda x, \lambda u, k/\lambda)$ leaves the dynamics unchanged and multiplies every cost by
$\lambda^2$, so the gap is the same. Starting twice as far from the origin is the same experiment as
doubling $k$.

\paragraph{Relation to the continuous-time construction.} The same idea is easier to state in
continuous time without discount. For $\dot x = f(x) + G(x)\, u$ with cost $q(x) + u^2$, the
Hamilton--Jacobi--Bellman equation gives $u^* = -\tfrac12 G^\top \nabla V$ and
\[
  \nabla V^\top f = -q + \tfrac14 \big(G^\top \nabla V\big)^2 ,
\]
which is linear in the drift $f$: choose $V$, $G$ and $q$, and any $f$ satisfying this one scalar
equation has $V$ as its cost-to-go. An example with $\beta(x_1) = \cos(2 x_1) + 2$ and
$q = x_1^2 + x_2^2$ is
\[
  \dot x_1 = -x_1 + x_2, \qquad
  \dot x_2 = -\tfrac12 x_1 - \tfrac12 x_2 \big(1 - \beta^2\big) + \beta\, u, \qquad
  V = \tfrac12 x_1^2 + x_2^2, \qquad u^* = -\beta\, x_2 ,
\]
whose $V$ is again the Riccati solution of the linearization, $P = \mathrm{diag}(1/2,\ 1)$. We do
not use it as the benchmark. The agents act in steps and with a discount, which is a different
optimization problem, and its optimal cost is not $V(x_0)$. From $x_0 = (1, 1)$, where $V = 1.5$,
the policy $u^*$ costs $1.68$ with Euler steps of $0.1$ and $1.52$ with steps of $0.01$; with a
discount of $0.99$ per step and steps of $0.01$ it costs $1.11$, below $V(x_0)$. A gap measured
against $V(x_0)$ would carry an error of that size, of either sign, as large as the floors in the
table. The construction of this section is the same idea carried out in discrete time, where the
Bellman equation is quadratic in the drift instead of linear; that is why a scale factor with a
square root appears in place of an added term.

%----------------------------------------------------------------------------------------------------
\section{A possible experiment}

\begin{itemize}[leftmargin=*]
  \item Train $A$, $B$ and $C$ on the scalar system with reward $r_t = -(x_t^2 + u_t^2)$, discount
    $\gamma = 0.99$ and linear policies $u = -kx$; a neural-network variant is evaluated by
    simulation. The baseline $K^*$, $J^*$ is computed with the same $\gamma$.
  \item Keep tuning seeds separate from evaluation seeds.
  \item Record $k_X(\omega, N)$ and $g_X(\omega, N)$ at checkpoints; gains outside
    $\mathcal{K}_\gamma$ score one, and the fraction of unstable gains, $|1 - k| \ge 1$, is reported.
  \item Report $\hat g_X$ with confidence bands, the floors, the bias and variance of the final gains,
    the crossing times, and which ordering of Section~\ref{sec:seeds} holds.
  \item Repeat on the nonlinear benchmark of Section~\ref{sec:nlqr} with $k = 2$, with the algorithms
    unchanged: they only see observations and rewards. Policies are scored by rollout, the gap is
    plotted on a logarithmic axis, and the floor of a linear policy ($0.0044$) is the reference
    for algorithms restricted to gains.
\end{itemize}

%----------------------------------------------------------------------------------------------------
\section{Limitations}
\label{sec:limitations}

\begin{itemize}[leftmargin=*]
  \item \textbf{Discounting.} $K^*$, $J^*$, $\mathcal{K}_\gamma$ and $c_\gamma$ depend on $\gamma$,
    so the baseline must use exactly the agents' $\gamma$. For $\gamma < 1$ finite cost does not
    imply stability, and in higher dimensions the discounted optimal gain itself need not stabilize
    the undiscounted system; stability is checked separately.
  \item \textbf{Noise.} $J$ becomes an expectation; in the scalar case the shape of $g$ is unchanged.
  \item \textbf{Linear policies.} They give the agent the form of the solution; the neural variant
    tests the realistic case.
  \item \textbf{Locality.} The quadratic approximation holds only near $K^*$, and the exact formula
    for the gap holds only for these parameter values.
  \item \textbf{Higher dimensions.} $K$ becomes a matrix and the expansion involves a Hessian.
  \item \textbf{The nonlinear benchmark is constructed.} Its drift is rescaled by $s(x)$ precisely so
    that the solution is exact; it is not a physical system, its nonlinearity is mild, and it has
    one particular structure (a state-dependent input gain). It tests whether an algorithm copes
    with a non-linear optimal policy, not nonlinear control in general.
  \item \textbf{SAC.} Evaluate it by its mean action to avoid the effect of the entropy bonus.
\end{itemize}

%----------------------------------------------------------------------------------------------------
\appendix

\section{Proximal Policy Optimization (PPO)}

PPO (Schulman et al.\ 2017) is an on-policy actor-critic algorithm. It alternates between collecting
data with the current policy and taking several gradient steps on a clipped objective that keeps the
new policy close to the one that collected the data.

\paragraph{Components.} A stochastic policy with parameters $\theta$ and a value function with
parameters $\phi$:
\[
  \pi_\theta(u \mid x), \qquad V_\phi(x) .
\]

\paragraph{Probability ratio.} How much more or less likely the new policy makes an action than the
policy that collected it:
\[
  \rho_t(\theta) = \frac{\pi_\theta(u_t \mid x_t)}{\pi_{\theta_{\mathrm{old}}}(u_t \mid x_t)} .
\]

\paragraph{Advantage estimate.} The temporal-difference errors of the value function, combined by
generalized advantage estimation (Schulman et al.\ 2016) with weight $\lambda$ and the discount
$\gamma$ of Section~\ref{sec:problem}, the same one used for the baseline:
\[
  \delta_t = r_t + \gamma V_\phi(x_{t+1}) - V_\phi(x_t), \qquad
  \hat A_t = \sum_{l=0}^{T - t - 1} (\gamma \lambda)^l\, \delta_{t+l}, \qquad
  \hat R_t = \hat A_t + V_\phi(x_t) .
\]

\paragraph{Clipped policy objective.} With clipping parameter $\varepsilon_{\mathrm{clip}}$:
\[
  L^{\mathrm{CLIP}}(\theta) = \hat{\mathbb{E}}_t\Big[\min\big(\rho_t(\theta)\, \hat A_t,\
  \operatorname{clip}\big(\rho_t(\theta),\, 1 - \varepsilon_{\mathrm{clip}},\,
  1 + \varepsilon_{\mathrm{clip}}\big)\, \hat A_t\big)\Big].
\]

\paragraph{Value loss and entropy bonus.}
\[
  L^{V}(\phi) = \hat{\mathbb{E}}_t\Big[\big(V_\phi(x_t) - \hat R_t\big)^2\Big], \qquad
  S[\pi_\theta](x_t) = -\int \pi_\theta(u \mid x_t) \log \pi_\theta(u \mid x_t)\, du .
\]

\paragraph{Combined objective.} Maximized with weights $c_1$ and $c_2$:
\[
  L(\theta, \phi) = \hat{\mathbb{E}}_t\Big[ L^{\mathrm{CLIP}}_t(\theta) - c_1 L^{V}_t(\phi)
  + c_2\, S[\pi_\theta](x_t) \Big].
\]

\paragraph{Algorithm.}
\begin{enumerate}[leftmargin=*]
  \item Initialize the policy parameters $\theta$ and the value parameters $\phi$.
  \item For each iteration:
    \begin{enumerate}
      \item Set the data-collecting policy to the current one, $\theta_{\mathrm{old}} \leftarrow \theta$.
      \item Run the policy for $T$ time steps (in one or more parallel environments) and store the
        states, controls and rewards.
      \item Compute the temporal-difference errors, the advantage estimates and the return targets
        from the stored data.
      \item For $E$ epochs, split the data into minibatches and, on each minibatch, take a gradient
        ascent step on the combined objective with respect to $\theta$ and $\phi$.
    \end{enumerate}
  \item Return the final policy.
\end{enumerate}

%----------------------------------------------------------------------------------------------------
\section*{References}

\begin{itemize}[leftmargin=*]
  \item Agarwal, R., Schwarzer, M., Castro, P.~S., Courville, A., and Bellemare, M.~G. (2021). Deep
    reinforcement learning at the edge of the statistical precipice. \emph{NeurIPS 2021}.
  \item Agyei, K., Sarhadi, P., and Polani, D. (2025).
    \href{https://arxiv.org/html/2507.08196}{Deep reinforcement learning in applied control:
    Challenges, analysis, and insights}. arXiv:2507.08196.
  \item Fazel, M., Ge, R., Kakade, S.~M., and Mesbahi, M. (2018).
    \href{https://proceedings.mlr.press/v80/fazel18a.html}{Global convergence of policy gradient
    methods for the linear quadratic regulator}. \emph{ICML 2018}.
  \item Henderson, P., Islam, R., Bachman, P., Pineau, J., Precup, D., and Meger, D. (2018). Deep
    reinforcement learning that matters. \emph{AAAI 2018}.
  \item Mania, H., Guy, A., and Recht, B. (2018).
    \href{https://arxiv.org/pdf/1803.07055}{Simple random search provides a competitive approach to
    reinforcement learning}. arXiv:1803.07055.
  \item Recht, B. (2019). \href{https://arxiv.org/pdf/1806.09460}{A tour of reinforcement learning:
    The view from continuous control}. \emph{Annual Review of Control, Robotics, and Autonomous
    Systems}. arXiv:1806.09460.
  \item Schulman, J., Moritz, P., Levine, S., Jordan, M., and Abbeel, P. (2016). High-dimensional
    continuous control using generalized advantage estimation. \emph{ICLR 2016}. arXiv:1506.02438.
  \item Schulman, J., Wolski, F., Dhariwal, P., Radford, A., and Klimov, O. (2017). Proximal policy
    optimization algorithms. arXiv:1707.06347.
  \item Tu, S., and Recht, B. (2019). \href{https://arxiv.org/pdf/1812.03565}{The gap between
    model-based and model-free methods on the linear quadratic regulator: An asymptotic viewpoint}.
    arXiv:1812.03565.
\end{itemize}

\end{document}
